\section{Introduction}
\label{sec:introduction}

Integrated assessment models (IAMs) link a model of the economy to a model of the climate and ask what a tonne of carbon dioxide costs society. Their answer, the social cost of carbon (SCC), is a shadow price: the marginal value of carbon in the atmosphere, relative to the marginal value of consumption. Computing it requires solving a dynamic optimisation problem over centuries, and the moment such a problem has more than a handful of state variables, or uncertainty, or tipping points, conventional grid-based dynamic programming runs into the curse of dimensionality. A growing literature therefore turns to methods from machine learning: deep equilibrium nets, which train a network to satisfy a model's equilibrium conditions \parencite{Azinovic2022DEQN,Maliar2021,Folini2024ClimateInCE}, their extension to uncertain parameters \parencite{friedl2023_deepUQ}, and reinforcement learning.

Operations research has a different and much older answer to the same problem. Stochastic dual dynamic programming (SDDP) \parencite{PereiraPinto1991} solves multistage stochastic programs, originally for hydrothermal power systems, by approximating each stage's value function from above with cutting planes \parencite{Kelley1960,Benders1962}. It alternates forward passes, which simulate the policy implied by the current cuts, with backward passes, which add a cut at each visited state using the value and the gradient of the stage problem. It needs convexity, but where convexity holds it comes with bounds on the optimum and convergence guarantees \parencite{Shapiro2011,Girardeau2015}. Neural SDDP \parencite{Dai2022NeuralSDDP} adds a network that maps a problem instance to its cuts, trained on earlier solutions, so that new instances can be solved faster or without iterations at all.

This paper applies both methods to the Analytic Climate Economy (ACE) of \textcite{TraegerChristian2023ACE}, in the DICE calibrations used by \textcite{traeger2021analytic_discussion}. ACE keeps the DICE production function, carbon cycle and temperature dynamics, but transforms capital and temperature so that the value function is affine and the model has a closed-form solution. That makes it a demanding benchmark: every number a method produces can be compared with the exact one.

The contributions are the following.

\begin{enumerate}
    \item \textbf{ACE as a sequence of convex programs.} With business-as-usual emissions on an exogenous path, which is the assumption behind ACE's closed form, each period's problem is convex as long as the cuts value capital positively (Proposition~\ref{prop:convex}). The gradient of the stage value, obtained from the multipliers on the cuts by the envelope theorem, is exactly the costate recursion of the model's equilibrium conditions (Proposition~\ref{prop:gradient}). The slopes of the cuts are therefore ACE's costates, and the SCC follows from the slopes of the next period's cut. This is a third route to the same shadow prices, next to the residuals of the equilibrium conditions used by deep equilibrium nets and the gradient of a learned value function used by actor-critic methods. Proposition~\ref{prop:convex}, the concavity of every value function on which the validity of the cuts rests, and the closed-form facts behind Proposition~\ref{prop:fixed} and Remark~\ref{rem:truncation} are verified formally in Lean (Appendix~\ref{app:lean}).
    \item \textbf{SDDP is exact in two iterations.} Because the value function is affine, a cut computed at any state is exact, the closed-form cut is a fixed point of the backward pass (Proposition~\ref{prop:fixed}), and SDDP from scratch converges after one backward pass (Corollary~\ref{cor:two}). Across 256 test draws of the time preference, the damage coefficient and the speed of atmosphere--ocean carbon exchange, it reproduces the closed-form SCC to within 0.002\%, in under three seconds per draw.
    \item \textbf{Neural SDDP predicts the cuts across parameter values.} A cut network trained on the SDDP solutions of 128 parameter draws gives, with a single forward pass and no SDDP iterations, a median SCC error of 0.24\% (DICE 2016) and 0.21\% (DICE 2023), at most 2.5\%, and a welfare loss below 0.0006\% on every test draw. On the same draws, this is 8 to 19 times more accurate than a pseudo-state deep equilibrium net trained over the same parameter box.
    \item \textbf{Warm starts do not help in ACE, and I explain why.} The predicted cuts make the first forward pass nearly optimal, within $10^{-9}$ of the optimum against 12\% below it from scratch. But certifying convergence requires a backward pass either way, and from scratch SDDP needs only one. The speed-ups of neural SDDP require value functions that need many cuts.
    \item \textbf{Uncertainty and the horizon.} With i.i.d.\ productivity shocks, SDDP proper (averaging cuts over shocks, sampling them forward) recovers the same costates and policy, as ACE predicts. The horizon after which the value is set to zero has to be long at low discount rates, because carbon's shadow price converges only at the rate $\beta$ (Remark~\ref{rem:truncation}).
\end{enumerate}

The paper is organised as follows. Section~\ref{sec:related} reviews related work. Section~\ref{sec:model} presents the model and its closed form. Section~\ref{sec:method} formulates the stage problems, derives the cuts and describes SDDP and neural SDDP. Section~\ref{sec:setup} describes the experiments, Section~\ref{sec:results} the results, and Section~\ref{sec:discussion} discusses what they imply for models without closed forms.
